\section{Konstruksi Dataset Penelitian}

Sebagaimana yag sudah disinggung sebelumnya, dataset ABAC dengan tingkat korelasi antar-atribut yang bervariasi dan berdomain IoT belum tersedia sebagai sumber daya siap pakai.
Penelusuran terhadap ABAC Lab \parencite{bui_abac_2025} menunjukkan bahwa lima dataset \textit{ground-truth} yang tersedia (\textit{healthcare, project-management, university, workforce, e-document}) seluruhnya berdomain organisasi/\textit{enterprise}, dan generator log sintetisnya belum mendukung kontrol eksplisit terhadap tingkat korelasi atribut.
Subbab ini menguraikan rancangan konstruksi dataset yang dibutuhkan untuk menguji hipotesis sebagai bagian dari aktivitas \textit{Design and Development}.

Konstruksi dataset dimulai dari dataset \textit{ground-truth} ABAC Lab beserta generator log sintetis \textit{open-source}-nya \parencite{bui_abac_2025} yang saat ini mendukung parameter ukuran log, rasio \textit{permit/deny}, dan ijeksi \textit{noise}.
Generator tersebut diperluas dengan menambahkan mekanisme injeksi korelasi atribut terkontrol.
Jadi, alih-alih men-\textit{sampling} nilai setiap atribut subjek dan objek secara independen, pasangan atribut yang relevan pada tiap domain (mis. departemen dan peran pengguna pada dataset \textit{workforce}) disampling dari distribusi bersama yang dikalibrasi untuk mencapai target \textit{Normalized Mutual Information} (NMI)\label{first:nmi} tertentu, sehingga tingkat korelasi yang di-\textit{inject} dapat dibandingkan langsung dengan struktur yang berhasil dipelajari algoritma.
Tiga tingkat korelasi ditetapkan sebagai variabel bebas (rendah, sedang, dan tinggi) dengan target NMI berturut-turut 0,05, 0,35, dan 0,65.
Batas bawah sengaja ditetapkan 0,05 dan bukan 0,00 karena entri berstatus \textit{permit} secara inheren berkorelasi sehingga memaksa proporsi tertentu pada log akan selalu menyuntikkan korelasi dasar yang tidak dapat dihilangkan.

Sebagai ilustrasi mengenai mekanisme penyuntikan korelasi, tinjau dua atribut pada domain \textit{workforce}, yaitu atribut subjek \texttt{position} dan atribut sumber daya \texttt{associatedContract}. 
Misalkan untuk suatu pasangan nilai tertentu berlaku probabilitas marginal
$P(\texttt{position}=p)=0{,}30$ dan
$P(\texttt{associatedContract}=c)=0{,}20$.
Apabila kedua atribut tersebut independen, probabilitas keduanya muncul secara bersamaan hanya sebesar
$0{,}30 \times 0{,}20 = 0{,}06$.
Dalam kondisi ini, mengetahui nilai \texttt{position} tidak memberikan informasi tambahan mengenai \texttt{associatedContract}, karena probabilitas bersyarat
$P(\texttt{associatedContract}=c \mid \texttt{position}=p)$
tetap sebesar 0,20. Pada proses penyuntikan korelasi, sebagian massa probabilitas dipindahkan menuju kombinasi nilai yang dibuat selaras tanpa mengubah distribusi marginal masing-masing atribut. 
Sebagai contoh ilustratif, apabila probabilitas gabungan pasangan tersebut dinaikkan dari 0,060 menjadi 0,151, maka
$P(\texttt{associatedContract}=c \mid \texttt{position}=p)$
meningkat dari 0,20 menjadi sekitar 0,503. 
Dengan demikian, frekuensi kemunculan masing-masing nilai atribut secara individual dapat tetap sama, tetapi pola kemunculan bersama keduanya menjadi tidak acak. 
Perubahan pada distribusi bersama inilah yang dimaksud sebagai peningkatan korelasi antar-atribut dalam konstruksi dataset penelitian ini.

\begin{table}[htbp]
    \centering
    \caption{Ilustrasi distribusi bersama dua atribut sebelum injeksi korelasi}
    \label{tab:ilustrasi-korelasi-sebelum}
    \begin{tabular}{lccc}
        \hline
         & \texttt{associatedContract} $=c$
         & \texttt{associatedContract} $\neq c$
         & Total \\
        \hline
        \texttt{position} $=p$
        & 0,060
        & 0,240
        & 0,300 \\
        \texttt{position} $\neq p$
        & 0,140
        & 0,560
        & 0,700 \\
        \hline
        Total
        & 0,200
        & 0,800
        & 1,000 \\
        \hline
    \end{tabular}
\end{table}

\begin{table}[htbp]
    \centering
    \caption{Ilustrasi distribusi bersama dua atribut setelah injeksi korelasi}
    \label{tab:ilustrasi-korelasi-sesudah}
    \begin{tabular}{lccc}
        \hline
         & \texttt{associatedContract} $=c$
         & \texttt{associatedContract} $\neq c$
         & Total \\
        \hline
        \texttt{position} $=p$
        & 0,151
        & 0,149
        & 0,300 \\
        \texttt{position} $\neq p$
        & 0,049
        & 0,651
        & 0,700 \\
        \hline
        Total
        & 0,200
        & 0,800
        & 1,000 \\
        \hline
    \end{tabular}
\end{table}

Tabel~\ref{tab:ilustrasi-korelasi-sebelum} dan
Tabel~\ref{tab:ilustrasi-korelasi-sesudah}
memperlihatkan bahwa penyuntikan korelasi tidak mengubah probabilitas marginal kedua atribut. 
Proporsi
\texttt{position} $=p$ tetap 0,30 dan proporsi
\texttt{associatedContract} $=c$ tetap 0,20.
Perubahan hanya terjadi pada distribusi bersama keduanya. 
Massa probabilitas pada kombinasi
\texttt{position} $=p$ dan
\texttt{associatedContract} $=c$
meningkat dari 0,060 menjadi 0,151, sedangkan massa pada sel silang berkurang secara bersesuaian. 
Akibatnya, probabilitas bersyarat berubah dari

\begin{equation}
P(\texttt{associatedContract}=c \mid \texttt{position}=p)
=
\frac{0{,}060}{0{,}300}
=
0{,}20
\end{equation}

menjadi

\begin{equation}
P(\texttt{associatedContract}=c \mid \texttt{position}=p)
=
\frac{0{,}151}{0{,}300}
\approx
0{,}503.
\end{equation}

Dengan kata lain, setelah korelasi disuntikkan, mengetahui nilai \texttt{position} memberikan informasi mengenai kemungkinan nilai \texttt{associatedContract}. 
Karakteristik inilah yang kemudian diukur melalui \textit{Normalized Mutual Information} (NMI) dan dapat dimanfaatkan oleh model yang mempertahankan dependensi antar-variabel.

Perbedaan tersebut penting bagi perbandingan algoritma dalam penelitian ini. 
Model univariat seperti \textit{compact-GA} dan \textit{compact-PSO} hanya mempertahankan probabilitas marginal masing-masing variabel sehingga tidak merepresentasikan secara eksplisit bahwa kemunculan suatu nilai \texttt{position} dapat meningkatkan kemungkinan kemunculan nilai \texttt{associatedContract} tertentu. 
Sebaliknya, CoDA-EDA mempelajari keterkaitan antar-variabel melalui \textit{mutual information} dan merepresentasikan dependensi terpilih sebagai sisi pada pohon Chow--Liu. 
Oleh karena itu, penyuntikan korelasi tidak sekadar mengubah frekuensi nilai atribut secara individual, tetapi membentuk struktur dependensi pada distribusi bersama yang dapat dimanfaatkan oleh model \textit{dependency-aware}. 
Nilai numerik pada ilustrasi ini hanya digunakan untuk memperjelas mekanisme. 
Pada dataset eksperimen, tingkat korelasi tidak ditetapkan langsung berdasarkan probabilitas gabungan pada ilustrasi tersebut, melainkan melalui proses kalibrasi \textit{closed-loop} hingga NMI empiris mendekati target 0,05, 0,35, dan 0,65.

Kalibrasi korelasi dilakukan secara \textit{closed-loop}.
Parameter penggandeng distribusi bersama dicari melalui pencarian biner hingga NMI yang terukur pada log final mendekati target, bukan hingga NMI distribusi teoretis mencapai target.
Prosedur ini diperlukan karena estimator NMI dan divergensi Jensen-Shannon sama-sama memiliki bias sampel-hingga yang membesar seiring rasio antara banyaknya sel kontingensi dan banyaknya sampel.
Pada pelaksanaan awal, kalibrasi yang memakai log rintisan berukuran kecil (800 entri) sementara log final berukuran 10.000 entri menghasilkan galat NMI sampai -0,11 pada dataset \textit{e-document} yang memiliki 379 sel kontingensi.
Galat tersebut menyusut menjadi -0,005 ketika ukuran log rintisan disamakan dengan ukuran log final.
Karena itu ukuran log kalibrasi ditetapkan sama dengan ukuran log final.

Bias serupa juga memengaruhi pengukuran \textit{drift} dan menuntut penyesuaian cara pelaporan.
Lantai derau divergensi Jensen-Shannon, yakni nilai yang terukur antar-jendela ketika distribusi sesungguhnya tidak berubah, bergantung pada rasio sel/sampel dan berbeda tajam antar-dataset.
Pada jendela rintisan berukuran kecil yang dipakai diagnostik awal, lantai derau terukur 0,017 pada \textit{workforce} (45 sel) tetapi mencapai 0,132-0,145 pada \textit{e-document} (379 sel).
Inilah yang memaksa pembesaran ukuran jendela.
Nilai ini bergantung ukuran jendela.
Pada ukuran jendela final yang lebih besar lantai derau \textit{e-document} menyusut ke kisaran 0,014-0,0333, namun tetap taktrivial (di atas $\tau = 0,01$) sehingga pelaporan \textit{drift} bersih tetap diperlukan.
Sifat sebagai properti struktural data ditegaskan oleh replikasinya yang konsisten pada dua lingkungan komputasi yang berbeda.
Konsekuensinya, target \textit{drift} 0,10 tidak mungkin tercapai dalam nilai mentah pada \textit{e-document}.
Tingkat \textit{drift} karena itu dilaporkan sebagai nilai bersih, yaitu divergensi terukur dikurangi lantai derau yang diperoleh dari kondisi stasioner.
Dengan koreksi ini, \textit{drift} terukur mendekati target pada kedua dataset (galat -0,006 hingga -0,022).
Implikasi metodologisnya, kondisi stasioner naik status.
Ia bukan sekadar kondisi pembanding dalam desain eksperimen, melainkan instrumen kalibrasi yang wajib dijalankan agar angka \textit{drift} dapat ditafsirkan.

Untuk menguji Hipotesis H3 dan H4 mengenai efisiensi dan kualitas pembaruan pada skenario \textit{drift}, log akses yang dihasilkan generator dibagi menjadi rangkaian jendela berurutan yang mensimulasikan siklus pembaruan kebijakan.
Antar-jendela, distribusi bersama atribut disesuaikan dengan besaran perubahan ($\Delta$) tertentu, yang diukur melalui divergensi Jensen-Shannon antara distribusi jendela berjalan dan jendela referensi sebelumnya, konsisten dengan definisi delta pada mekanisme deteksi \textit{drift} CoDA-EDA.

Sebelum digunakan pada aktivitas \textit{Demonstration}, setiap varian dataset yang dikonstruksi akan diverifikasi kesesuaiannya dengan target rancangan: tingkat korelasi aktual (diukur via NMI pada data yang dihasilkan) dibandingkan terhadap target yang ditetapkan, dan proporsi keputusan \textit{permit/deny} diperiksa agar tetap berada pada rentang wajar (tidak timpang secara ekstrem) sebagaimana disyaratkan oleh parameter rasio \textit{permit/deny} bawaan ABAC Lab.
Keseluruhan alur konstruksi dataset, dari sumber dasar hingga dataset penelitian akhir, dirangkum pada Gambar \ref{fig:konstruksi_dataset}.

Pelaksanaan awal konstruksi dataset pada perangkat target menghasilkan tiga temuan yang mengubah rancangan pada subbab ini, dan karenanya dicatat secara eksplisit demi keterulangan.
Temuan pertama menyangkut kelangkaan \textit{permit}.
Enumerasi penuh ruang $(subjek \times objek \times aksi)$ menunjukkan bahwa proporsi permintaan berstatus \textit{permit} hanya berkisar 2-5\% pada seluruh dataset ABAC Lab (\textit{workforce} 1,99\%; \textit{university} 2,50\%; \textit{project-management} 3,32\%; \textit{healthcare} 4,27\%; \textit{e-document} 5,34\%).
Konsekuensinya, log tidak dapat dibangkitkan melalui sampling seragam atas ruang permintaan, melainkan harus melalui stratifikasi \textit{permit/deny} dengan himpunan \textit{permit} dipra-komputasi terlebih dahulu.

Temuan kedua menyangkut kelayakan dataset.
Seluruh dataset ABAC Lab yang digunakan bersifat sintetis yang dibangkitkan oleh generator log ABAC Lab \parencite{bui_abac_2025} yang diperluas sehingga karakteristik yang dilaporkan berikut merupakan properti konstruksi, bukan properti data operasional.
Setelah dibatasi pada entitas yang memiliki pasangan atribut terkorelasi, jumlah \textit{triple permit} unik yang tersedia berbeda tajam antar-dataset (\textit{healthcare} hanya 14, \textit{project-management} 17, dan \textit{university} 90, dibandingkan \textit{workforce} 14.513 dan \textit{e-document} 32.961).
Pada ketiga dataset kecil, pembangkitkan log berukuran memadai menghasilkan tingkat duplikasi entri 48-93\%, sehingga tugas \textit{mining} menjadi \textit{degenerate}, yakni algoritma dihadapkan pada pengulangan pola yang sama alih-alih keragaman pola yang wajar.
Karena itu penelitian ini menetapkan \textit{workforce} dan \textit{e-document} sebagai dataset utama untuk pengujian Hipotesis H1, sebab \textit{permit} uniknya memberi keragaman pola yang memadai.
Ketiga dataset kecil tetap disertakan pada pengujian seluruh hipotesis, namun dengan peran yang disesuaikan menurut karakteristiknya.
Untuk H1, hasil pada ketiganya diperlakukan sebagai konteks pelengkap dan bukan tumpuan vonis, karena tingginya duplikasi entri menumpulkan pembeda antar-algoritma.
Sedangkan untuk H2 serta H3 dan H4, ketiganya sepenuhnya sahih, bahkan lantai derau divergensi informasinya yang rendah menjadikannya lingkungan yang justru layak untuk menunjukkan mekanisme H3.
Dimensi ruang solusi, yaitu banyaknya pasangan atribut-nilai yang membentuk vektor keputusan adalah D = 436 untuk \textit{workforce} dan D = 1.907 untuk \textit{e-document}, serta D = 46, 51, dan 54 untuk \textit{healthcare, project-management}, dan \textit{university}.

Temuan ketiga menyangkut hubungan antara korelasi dan duplikasi.
Tingkat duplikasi entri meningkat seiring naiknya target korelasi (\textit{workforce} $17\% \to 37\% \to 51\%$; \textit{e-document} $4\% \to 14\% \to 30\%$ untuk korelasi rendah $\to$ sedang $\to$ tinggi).
Peningkatan ini bersifat struktural, bukan cacat implementasi.
Korelasi tinggi secara matematis memusatkan sampling pada sedikit sel distribusi bersama, sehingga entri yang sama lebih sering terpilih.
Duplikasi semacam ini juga tidak menyalahi realisme, karena log akses nyata memang memuat permintaan berulang.
Meskipun demikian, tingkat duplikasi tetap dilaporkan pada setiap konfigurasi agar pembaca dapat menilai keragaman efektif data yang dihadapi algoritma.

\begin{figure}[H]
    \centering
    \includegraphics[width=0.9\textwidth]{gambar_dataset.jpg}
    \caption{\textit{Pipeline} Konstruksi Dataset Penelitian}
    \label{fig:konstruksi_dataset}
\end{figure}

Gambar \ref{fig:konstruksi_dataset} menegaskan satu keputusan rancang yang penting, yakni metrik yang dipakai untuk menyuntikkan korelasi dan \textit{drift} ke dataset sengaja dipilih identik dengan metrik yang dipakai CoDA-EDA untuk mempelajari struktur dependensi dan mendeteksi \textit{drift}.
Dengan demikian, kadar korelasi dan \textit{drift} yang tertanam pada dataset dapat dibandingkan langsung dengan yang berhasil ditangkap algoritma, sehingga evaluasi H1, H3, dan H4 memiliki acuan \textit{ground-truth} yang jelas.

Di luar dataset yang dikonstruksi sendiri, penelitian ini turut memakai dataset \textit{Smart Building IoT} \parencite{diaz-rodriguez_abac_2025} sebagai validasi eksternal, yakni memeriksa apakah temuan bertahan pada data yang dibangkitkan pihak lain dengan generator, kebijakan, dan distribusi yang berbeda.
Dataset ini digunakan apa adanya tanpa injeksi korelasi atau \textit{drift} buatan.
Karena berukuran 5,7 juta permintaan akses dengan rasio \textit{permit} 0,75 dan dimensi penuh sekitar 1.164, ia disubsampel secara terstratifikasi (mempertahankan rasio \textit{permit} asli) dan nilai atribut ber-\textit{support} rendah disatukan, menghasilkan dimensi efektif sekitar 428.
Perlu ditegaskan bahwa dataset tersebut bersifat sintetis sebagaimana dinyatakan penulisnya.
Nilai validasinya terletak pada asal generatornya yang independen, bukan pada keasliannya sebagai data operasional.

\begin{longtable}{|c|>{\raggedright\arraybackslash}p{2.2cm}|p{1.5cm}|p{1.5cm}|>{\raggedright\arraybackslash}p{3cm}|>{\raggedright\arraybackslash}p{3cm}|}
\caption{Skema Dataset Penelitian}
\label{tab:skema-dataset} \\
\hline
\textbf{No} & \textbf{Dataset Dasar} & \textbf{Dimensi $D$} & \textbf{Permit Unik} & \textbf{Perluasan yang Diterapkan} & \textbf{Peran dalam Penelitian} \\
\hline
\endfirsthead

\multicolumn{6}{c}{{\bfseries \tablename\ \thetable{} -- lanjutan dari halaman sebelumnya}} \\
\hline
\textbf{No} & \textbf{Dataset Dasar} & \textbf{Dimensi $D$} & \textbf{Permit Unik} & \textbf{Perluasan yang Diterapkan} & \textbf{Peran dalam Penelitian} \\
\hline
\endhead

\hline
\multicolumn{6}{r}{Lanjut ke halaman berikutnya} \\
\endfoot

\hline
\endlastfoot

% ==================== DATA BARIS ====================
1 & \textit{Workforce} (ABAC Lab) & 436 & 14.513 & 3 tingkat korelasi (NMI rendah/sedang/tinggi) $\times$ 3 skenario drift (stasioner/rendah/sedang) & EKSPERIMEN UTAMA --- pengujian H1 (korelasi) dan H3-H4 (drift) \\
\hline
2 & \textit{E-document} (ABAC Lab) & 1.907 & 32.961 & 3 tingkat korelasi $\times$ 3 skenario drift & EKSPERIMEN UTAMA --- pengujian H1 dan H3-H4 pada dimensi lebih besar \\
\hline
3 & \textit{University} (ABAC Lab) & 54 & 90 & 3 tingkat korelasi $\times$ 3 skenario drift & KONTEKS PELENGKAP --- H2/H3/H4 sahih; vonis H1 terbatas (duplikasi 27-57\% pada log memadai) \\
\hline
4 & \textit{Project-management} (ABAC Lab) & 51 & 17 & 3 tingkat korelasi $\times$ 3 skenario drift & KONTEKS PELENGKAP --- H2/H3/H4 sahih; vonis H1 terbatas (duplikasi 48-69\%) \\
\hline
5 & \textit{Healthcare} (ABAC Lab) & 46 & 14 & 3 tingkat korelasi $\times$ 3 skenario drift & KONTEKS PELENGKAP --- H2/H3/H4 sahih; vonis H1 terbatas (duplikasi 58-93\%) \\
\hline
6 & \textit{Smart Building IoT} \parencite{diaz-rodriguez_abac_2025} & 428* & --- & Digunakan apa adanya; subsampel terstratifikasi dari 5,7 juta baris, tanpa injeksi korelasi/drift & VALIDASI EKSTERNAL --- generalisasi H1-H4 pada dataset berdomain IoT yang dibangkitkan pihak lain \\
\hline

\end{longtable}

Dataset-dataset pada Tabel \ref{tab:skema-dataset} menjadi lingkungan pengujian bersama bagi CoDA-EDA dan seluruh algoritma pembanding.
Tanda (*) pada dataset \textit{Smart Building IoT} menunjukkan dimensi efektif setelah subsampel dan penyatuan nilai ber\textit{support} rendah yang mana berkas asalnya memuat 5,7 juta permintaan akses dengan rasio \textit{permit} 0,75.